\chapter{जीन से प्रोटीन तक}\label{ch:g11:gene-expression}

1961 में किसी जीव-रसायनज्ञ ने जीवाणुओं के कोशिका-रहित सत्त को ऐसा कृत्रिम
\omterm{def:g11:gene-expression:transcription}{RNA} खिलाया जो केवल U अक्षर के दोहराव से बना था। उस सत्त ने ऐसा \omterm{def:g11:gene-expression:protein}{प्रोटीन}
बनाया जो केवल एक ही \omterm{def:g10:chemistry-of-life:families}{ऐमीनो अम्ल}, फेनिलऐलानिन, के दोहराव से बना था। उस
अकेले प्रयोग से \omterm{def:g11:gene-expression:code}{आनुवंशिक कूट} का पहला शब्द पढ़ लिया गया: UUU का अर्थ है
फेनिलऐलानिन। पाँच साल के भीतर चौंसठों शब्द ज्ञात हो गए, और वे किसी
जीवाणु, गेहूँ के पौधे और मनुष्य में वही शब्द निकले। यह अध्याय किसी \omterm{def:g10:universal-dna:gene}{जीन} के
\omterm{def:g10:universal-dna:nucleotide}{न्यूक्लियोटाइड} अनुक्रम से किसी \omterm{def:g11:gene-expression:protein}{प्रोटीन} के \omterm{def:g10:chemistry-of-life:families}{ऐमीनो अम्ल} अनुक्रम तक के रास्ते
पर चलता है --- वही रास्ता जिस पर \cref{ch:g11:mutations} का हर लक्षण
दौड़ता है।

\section{जीन प्रोटीन बनाते हैं}

\begin{proposition}[एक जीन, एक प्रोटीन]\label{prop:g11:gene-expression:onegene}
कोई \omterm{def:g10:universal-dna:gene}{जीन} तब अभिव्यक्त होता है जब \omterm{def:g10:cells-common-unit:cell}{कोशिका} उसके अनुक्रम से वह \omterm{def:g11:gene-expression:protein}{प्रोटीन} बनाती है
जिसका वह कूट है। अधिकांश लक्षण प्रोटीनों पर टिके हैं --- \omterm{def:g10:cell-metabolism:metabolism}{एंजाइम}, संरचनात्मक
तंतु, वाहक, ग्राही --- और कोई \omterm{def:g11:mutations:mutation}{उत्परिवर्तन} किसी लक्षण को उसके \omterm{def:g10:universal-dna:gene}{जीन} के बनाए
\omterm{def:g11:gene-expression:protein}{प्रोटीन} को बदलकर ही बदलता है।
\end{proposition}
\begin{proof}[साक्ष्य]
बीडल और टैटम (1941) ने ब्रेड की फफूँद के बीजाणुओं पर विकिरण डाला और ऐसे
उत्परिवर्ती इकट्ठे किए जो न्यूनतम माध्यम पर तब तक नहीं उग सकते थे जब तक
कोई एक ख़ास पदार्थ, जैसे \omterm{def:g10:chemistry-of-life:families}{ऐमीनो अम्ल} आर्जिनीन, मिलाया न जाए। हर उत्परिवर्ती
में अभिक्रियाओं की उस शृंखला का, जो वह पदार्थ बनाती है, कोई एक \omterm{def:g10:cell-metabolism:metabolism}{एंजाइम} नहीं
था; हर ख़राबी अकेले एक \omterm{def:g10:universal-dna:gene}{जीन} की तरह विरासत में मिलती थी; और उसी शृंखला के
अलग-अलग क़दमों पर रुके उत्परिवर्ती अलग-अलग \omterm{def:g10:universal-dna:gene}{जीनों} में \omterm{def:g11:mutations:mutation}{उत्परिवर्तन} ढोते थे।
एक \omterm{def:g10:universal-dna:gene}{जीन}, एक \omterm{def:g10:cell-metabolism:metabolism}{एंजाइम} --- और, जैसा बाद के काम ने सामान्य किया, एक \omterm{def:g10:universal-dna:gene}{जीन}, एक
पॉलिपेप्टाइड शृंखला।
\end{proof}

\begin{definition}[प्रोटीन, ऐमीनो अम्ल अनुक्रम]\label{def:g11:gene-expression:protein}
\emph{प्रोटीन}\index{प्रोटीन} \omterm{def:g10:chemistry-of-life:families}{ऐमीनो अम्लों} की एक शृंखला है, जो कुछ दर्जन
से लेकर कई हज़ार तक लंबी होती है और बीस क़िस्मों के समुच्चय से ली जाती है।
\emph{अनुक्रम} --- यानी \omterm{def:g10:chemistry-of-life:families}{ऐमीनो अम्लों} का क्रम --- ही वह चीज़ है जिसे \omterm{def:g10:universal-dna:gene}{जीन} तय
करता है; और बन जाने के बाद वह शृंखला उसी अनुक्रम से तय एक निश्चित
त्रि-आयामी आकार में मुड़ जाती है, तथा वही आकार उसका काम तय करता है। एक
\omterm{def:g10:chemistry-of-life:families}{ऐमीनो अम्ल} बदलिए और वह मोड़, तथा वह काम, बदल सकता है।
\end{definition}

\section{अनुलेखन: जीन की RNA में नक़ल}

\begin{definition}[संदेशवाहक RNA और अनुलेखन]\label{def:g11:gene-expression:transcription}
\emph{RNA}\index{RNA} (राइबोन्यूक्लिक अम्ल) \omterm{def:g10:universal-dna:information}{DNA} जैसे ही \omterm{def:g10:universal-dna:nucleotide}{न्यूक्लियोटाइडों}
की इकहरी लड़ी वाली शृंखला है, बस उसकी शर्करा राइबोस है और थाइमीन (T) की
जगह \omterm{def:g10:universal-dna:nucleotide}{क्षारक} यूरैसिल (U) आता है। \emph{अनुलेखन}\index{अनुलेखन} किसी \omterm{def:g10:universal-dna:gene}{जीन} के
अनुक्रम की \emph{संदेशवाहक RNA}\index{संदेशवाहक RNA} (mRNA) के एक अणु में
नक़ल है: \omterm{def:g10:cell-metabolism:metabolism}{एंजाइम} \emph{RNA पॉलिमरेज़} \omterm{def:g10:universal-dna:gene}{जीन} के साथ-साथ दोहरी कुंडली खोलता
है, एक लड़ी (यानी \emph{साँचा लड़ी}) पढ़ता है और उसके पूरक RNA को जोड़ देता
है, जिसमें U के सामने A, A के सामने T, G के सामने C और C के सामने G आता
है। उस mRNA का अनुक्रम \omterm{def:g10:universal-dna:gene}{जीन} की दूसरी लड़ी वाला होता है, बस T की जगह U; और
वह \omterm{def:g10:cells-common-unit:organelle}{केंद्रक} छोड़कर \omterm{def:g10:cells-common-unit:cell}{कोशिकाद्रव्य} में चला जाता है।
\end{definition}

\begin{omfigure}
\begin{tikzpicture}[font=\small, >=Stealth, line cap=round]
  % DNA strands with bubble
  \draw[very thick, omDef] (0,0.5) -- (2.5,0.5) .. controls (3.2,0.5) and (3.4,1.3) .. (4.3,1.3) -- (6.7,1.3) .. controls (7.6,1.3) and (7.8,0.5) .. (8.5,0.5) -- (11,0.5);
  \draw[very thick, omDef] (0,-0.5) -- (2.5,-0.5) .. controls (3.2,-0.5) and (3.4,-1.3) .. (4.3,-1.3) -- (6.7,-1.3) .. controls (7.6,-1.3) and (7.8,-0.5) .. (8.5,-0.5) -- (11,-0.5);
  \foreach \x in {0.3,0.7,...,2.3} \draw[gray] (\x,0.5) -- (\x,-0.5);
  \foreach \x in {8.7,9.1,...,10.9} \draw[gray] (\x,0.5) -- (\x,-0.5);
  % template strand letters (bottom strand in bubble)
  \foreach \x/\l in {4.5/T, 4.9/A, 5.3/C, 5.7/G, 6.1/G, 6.5/A} \node[font=\scriptsize, omDef, anchor=north] at (\x,-1.32) {\l};
  % mRNA
  \draw[very thick, omThm] (4.4,-0.9) -- (6.6,-0.9) .. controls (7.4,-0.9) and (7.2,0.1) .. (8.0,0.6) .. controls (8.6,1.0) and (9.4,1.4) .. (10.6,1.7);
  \foreach \x/\l in {4.5/A, 4.9/U, 5.3/G, 5.7/C, 6.1/C, 6.5/U} \node[font=\scriptsize, omThm, anchor=south] at (\x,-0.85) {\l};
  \foreach \x in {4.5,4.9,...,6.5} \draw[gray, densely dotted] (\x,-0.95) -- (\x,-1.28);
  % polymerase
  \draw[thick, fill=omProp!30, rounded corners=8pt] (6.95,-1.6) rectangle (8.65,0.2);
  \node[align=center] at (7.8,-0.7) {RNA\\ पॉलिमरेज़};
  \draw[->, thick] (8.7,-1.9) -- (9.8,-1.9) node[right] {जीन के साथ-साथ चलता है};
  \node[anchor=east, omDef, align=right] at (-0.15,0) {जीन का\\ DNA};
  \node[anchor=west, omThm, align=left] at (10.7,1.7) {mRNA, जो छूट\\ रहा है};
  \node[anchor=north, omDef] at (5.5,-1.75) {साँचा लड़ी};
\end{tikzpicture}

{\small \omterm{def:g11:gene-expression:transcription}{अनुलेखन}। \omterm{def:g11:gene-expression:transcription}{RNA} पॉलिमरेज़ कुंडली खोलता है, साँचा लड़ी पढ़ता है और
उसका पूरक \omterm{def:g11:gene-expression:transcription}{संदेशवाहक RNA} बनाता है (T की जगह U)। \omterm{def:g10:cell-metabolism:metabolism}{एंजाइम} के पीछे कुंडली
दोबारा बंद हो जाती है; और वह \omterm{def:g11:gene-expression:transcription}{RNA} छिलकर अलग हो जाता है।}
\end{omfigure}

\begin{example}[एक अनुलेख]\label{ex:g11:gene-expression:transcript}
साँचा लड़ी \texttt{3'-TAC GGA CTT ATC-5'} से mRNA
\texttt{5'-AUG CCU GAA UAG-3'} बनता है। \omterm{def:g10:universal-dna:gene}{जीन} की दूसरी लड़ी
\texttt{5'-ATG CCT GAA TAG-3'} पढ़ी जाती है: mRNA उसी लड़ी का अनुक्रम है,
बस T की जगह U के साथ, और इसीलिए जीन-अनुक्रम परिपाटी के अनुसार उसी लड़ी के
रूप में लिखे जाते हैं, यानी \emph{कूट लड़ी} के रूप में।
\end{example}

\section{आनुवंशिक कूट}

\begin{definition}[कोडॉन और आनुवंशिक कूट]\label{def:g11:gene-expression:code}
mRNA किसी तय शुरुआती बिंदु से तीन-तीन \omterm{def:g10:universal-dna:nucleotide}{न्यूक्लियोटाइडों} के लगातार समूहों
में, यानी \emph{कोडॉनों}\index{कोडॉन} में, पढ़ा जाता है।
\emph{आनुवंशिक कूट}\index{आनुवंशिक कूट} 64 संभव कोडॉनों और 20 \omterm{def:g10:chemistry-of-life:families}{ऐमीनो
अम्लों} के बीच का मेल है: 61 कोडॉन कोई \omterm{def:g10:chemistry-of-life:families}{ऐमीनो अम्ल} तय करते हैं, तीन (UAA,
UAG, UGA) \emph{विराम} संकेत हैं जो शृंखला ख़त्म कर देते हैं, और AUG, जो
मेथिओनीन तय करता है, \emph{आरंभ} कोडॉन का भी काम करता है। यह कूट
\emph{अपह्रासित} है --- अधिकांश \omterm{def:g10:chemistry-of-life:families}{ऐमीनो अम्लों} के कई कोडॉन हैं --- और
\emph{सार्वत्रिक}: वही तालिका हर ज्ञात जीव के काम आती है, विरले छोटे-मोटे
भेदों के साथ।
\end{definition}

\begin{omfigure}
\begin{tabular}{c|cccc|c}
\multicolumn{1}{c}{} & \multicolumn{4}{c}{\small दूसरा अक्षर} & \multicolumn{1}{c}{} \\
{\small पहला} & U & C & A & G & {\small तीसरा} \\ \hline
U & UUU Phe & UCU Ser & UAU Tyr & UGU Cys & U \\
 & UUC Phe & UCC Ser & UAC Tyr & UGC Cys & C \\
 & UUA Leu & UCA Ser & UAA \textbf{विराम} & UGA \textbf{विराम} & A \\
 & UUG Leu & UCG Ser & UAG \textbf{विराम} & UGG Trp & G \\ \hline
C & CUU Leu & CCU Pro & CAU His & CGU Arg & U \\
 & CUC Leu & CCC Pro & CAC His & CGC Arg & C \\
 & CUA Leu & CCA Pro & CAA Gln & CGA Arg & A \\
 & CUG Leu & CCG Pro & CAG Gln & CGG Arg & G \\ \hline
A & AUU Ile & ACU Thr & AAU Asn & AGU Ser & U \\
 & AUC Ile & ACC Thr & AAC Asn & AGC Ser & C \\
 & AUA Ile & ACA Thr & AAA Lys & AGA Arg & A \\
 & AUG \textbf{Met} & ACG Thr & AAG Lys & AGG Arg & G \\ \hline
G & GUU Val & GCU Ala & GAU Asp & GGU Gly & U \\
 & GUC Val & GCC Ala & GAC Asp & GGC Gly & C \\
 & GUA Val & GCA Ala & GAA Glu & GGA Gly & A \\
 & GUG Val & GCG Ala & GAG Glu & GGG Gly & G \\
\end{tabular}

{\small \omterm{def:g11:gene-expression:code}{आनुवंशिक कूट}, जो mRNA पर पढ़ा जाता है। हर \omterm{def:g11:gene-expression:code}{कोडॉन} अपने पहले अक्षर
(पंक्ति), दूसरे अक्षर (स्तंभ) और तीसरे अक्षर (खंड के भीतर की पंक्ति) से
मिलता है। AUG मेथिओनीन भी है और आरंभ का संकेत भी; और तीन \omterm{def:g11:gene-expression:code}{कोडॉन} विराम हैं।}
\end{omfigure}

\begin{proposition}[कूट पढ़ा कैसे गया]\label{prop:g11:gene-expression:readcode}
हर \omterm{def:g11:gene-expression:code}{कोडॉन} का अर्थ प्रयोग से स्थापित किया गया था।
\end{proposition}
\begin{proof}[साक्ष्य]
नीरेनबर्ग और मात्थाई (1961) ने \omterm{def:g10:cells-common-unit:organelle}{राइबोसोम}, \omterm{def:g10:chemistry-of-life:families}{ऐमीनो अम्ल} और प्रोटीन-संश्लेषण
के \omterm{def:g10:cell-metabolism:metabolism}{एंजाइम} रखते किसी कोशिका-रहित सत्त में अकेले एक दोहराए गए
\omterm{def:g10:universal-dna:nucleotide}{न्यूक्लियोटाइड} का कृत्रिम \omterm{def:g11:gene-expression:transcription}{RNA} मिलाया: पॉली-U ने फेनिलऐलानिनों की शृंखला
दी, पॉली-C ने प्रोलीनों की, और पॉली-A ने लाइसीनों की। दोहराए गए जोड़ों और
तिकड़ियों वाले कृत्रिम \omterm{def:g11:gene-expression:transcription}{RNA} ने, और बाद में इकहरी तिकड़ियों के \omterm{def:g10:cells-common-unit:organelle}{राइबोसोमों} से
बँधने ने, बचे हुए \omterm{def:g11:gene-expression:code}{कोडॉन} 1966 तक तय कर दिए। तीन अक्षरों की लंबाई गणित से
अनुमानित की गई थी (दो अक्षर केवल 16 मेल देते हैं, जो 20 \omterm{def:g10:chemistry-of-life:families}{ऐमीनो अम्लों} के
लिए बहुत कम हैं; तीन 64 देते हैं) और \omterm{def:g11:mutations:mutation}{उत्परिवर्तनों} ने उसकी पुष्टि की: किसी
\omterm{def:g10:universal-dna:gene}{जीन} में एक या दो \omterm{def:g10:universal-dna:nucleotide}{न्यूक्लियोटाइड} डालने से उसका \omterm{def:g11:gene-expression:protein}{प्रोटीन} नष्ट हो जाता है, और
तीन डालने से लगभग सामान्य \omterm{def:g11:gene-expression:protein}{प्रोटीन} लौट आता है।
\end{proof}

\begin{method}[किसी अनुक्रम का अनुवादन]\label{met:g11:gene-expression:translate}
\begin{enumerate}
  \item कूट लड़ी लिखिए (या साँचा लड़ी का \omterm{def:g11:gene-expression:transcription}{अनुलेखन} कीजिए): mRNA पाने के लिए
        T की जगह U रख दीजिए।
  \item पहला AUG ढूँढ़िए: वही आरंभ है, और वाचन चौखट उसी से तय हो जाती है।
  \item उस AUG से mRNA को लगातार तिकड़ियों में काटिए और हर एक को तालिका
        में पढ़िए।
  \item पहले विराम \omterm{def:g11:gene-expression:code}{कोडॉन} पर रुक जाइए; क्रम से गिनाए गए वे \omterm{def:g10:chemistry-of-life:families}{ऐमीनो अम्ल} ही
        \omterm{def:g11:gene-expression:protein}{प्रोटीन} का अनुक्रम हैं।
  \item किसी उत्परिवर्ती के लिए यही दोहराइए और तुलना कीजिए: वही \omterm{def:g11:gene-expression:protein}{प्रोटीन}
        (मूक), एक \omterm{def:g10:chemistry-of-life:families}{ऐमीनो अम्ल} बदला (कुअर्थी), समय से पहले विराम (निरर्थक),
        \omterm{def:g11:mutations:mutation}{उत्परिवर्तन} के बाद सब कुछ बदला (चौखट-खिसकाव)।
\end{enumerate}
\end{method}

\begin{example}[चार उत्परिवर्तन, चार नतीजे]\label{ex:g11:gene-expression:outcomes}
कूट लड़ी \texttt{ATG CCT GAA TTC TAG}: mRNA
\texttt{AUG CCU GAA UUC UAG}, \omterm{def:g11:gene-expression:protein}{प्रोटीन} Met--Pro--Glu--Phe।
\begin{itemize}
  \item \texttt{ATG CC\textbf{C} GAA TTC TAG}: CCC अब भी Pro है ---
        यानी \emph{मूक}।
  \item \texttt{ATG CCT G\textbf{C}A TTC TAG}: GCA यानी Ala ---
        Met--Pro--Ala--Phe, यानी \emph{कुअर्थी}।
  \item \texttt{ATG CCT \textbf{T}AA TTC TAG}: UAA यानी विराम ---
        Met--Pro, यानी \emph{निरर्थक}, एक कटा हुआ \omterm{def:g11:gene-expression:protein}{प्रोटीन}।
  \item \texttt{ATG CC\textbf{A}T GAA TTC TAG}: डाला गया एक A चौखट खिसका
        देता है: \texttt{AUG CCA UGA}\dots{} Met--Pro--विराम ---
        यानी \emph{चौखट-खिसकाव}।
\end{itemize}
\end{example}

\section{अनुवादन: संदेश पढ़ा गया}

\begin{proposition}[अनुवादन]\label{prop:g11:gene-expression:translation}
\emph{अनुवादन}\index{अनुवादन} \omterm{def:g10:cells-common-unit:cell}{कोशिकाद्रव्य} में \emph{\omterm{def:g10:cells-common-unit:organelle}{राइबोसोमों}} पर होता
है। कोई \omterm{def:g10:cells-common-unit:organelle}{राइबोसोम} mRNA से उसके आरंभ \omterm{def:g11:gene-expression:code}{कोडॉन} पर बँधता है और उसके साथ-साथ
कोडॉन-दर-कोडॉन चलता है। हर \omterm{def:g11:gene-expression:code}{कोडॉन} से कोई \emph{अंतरण \omterm{def:g11:gene-expression:transcription}{RNA}} (tRNA) मेल खाता
है, यानी छोटा-सा ऐसा \omterm{def:g11:gene-expression:transcription}{RNA} जो एक सिरे पर उस \omterm{def:g11:gene-expression:code}{कोडॉन} का पूरक तीन-न्यूक्लियोटाइड
वाला \emph{प्रतिकोडॉन} ढोता है और दूसरे सिरे पर उससे मेल खाता \omterm{def:g10:chemistry-of-life:families}{ऐमीनो अम्ल}।
\omterm{def:g10:cells-common-unit:organelle}{राइबोसोम} हर लाए गए \omterm{def:g10:chemistry-of-life:families}{ऐमीनो अम्ल} को बढ़ती शृंखला से जोड़ देता है और ख़ाली
tRNA छोड़ देता है; और किसी विराम \omterm{def:g11:gene-expression:code}{कोडॉन} पर वह पूरी हुई शृंखला छोड़ देता है,
जो मुड़ जाती है। कई \omterm{def:g10:cells-common-unit:organelle}{राइबोसोम} एक ही mRNA को एक के बाद एक पढ़ते हैं, और एक
mRNA नष्ट होने से पहले उस \omterm{def:g11:gene-expression:protein}{प्रोटीन} की सैकड़ों प्रतियाँ दे देता है।
\end{proposition}
\admitted

\begin{omfigure}
\begin{tikzpicture}[font=\small, >=Stealth, line cap=round]
  % mRNA
  \draw[very thick, omThm] (0,0) -- (11,0);
  \foreach \x/\l in {0.4/A,0.8/U,1.2/G, 1.7/C,2.1/C,2.5/U, 3.0/G,3.4/A,3.8/A, 4.3/U,4.7/U,5.1/C, 5.6/G,6.0/G,6.4/U, 6.9/U,7.3/A,7.7/G}
    \node[font=\scriptsize, omThm, anchor=north] at (\x,-0.05) {\l};
  \node[anchor=east, omThm] at (-0.15,0) {mRNA};
  % ribosome (two lobes) over codons 3 and 4
  \draw[thick, fill=omProp!25, rounded corners=10pt] (2.7,0.15) rectangle (5.5,1.4);
  \draw[thick, fill=omProp!35, rounded corners=12pt] (2.5,1.4) .. controls (2.5,3.0) and (5.7,3.0) .. (5.7,1.4) -- cycle;
  \node at (4.1,2.1) {राइबोसोम};
  % tRNAs
  \draw[very thick, omMeth!80!black] (3.4,0.55) -- (3.4,1.35);
  \node[font=\scriptsize, omMeth!80!black, anchor=south] at (3.4,0.15) {CUU};
  \draw[very thick, omMeth!80!black] (4.7,0.55) -- (4.7,1.35);
  \node[font=\scriptsize, omMeth!80!black, anchor=south] at (4.7,0.15) {AAG};
  % amino acids chain
  \foreach \i/\x in {0/3.4, 1/4.7} \fill[omDef!70] (\x,1.5) circle (4pt);
  \foreach \x in {2.6, 1.9, 1.2} \fill[omDef!70] (\x,1.9) circle (4pt);
  \draw[thick, omDef!70] (1.2,1.9) -- (2.6,1.9) -- (3.4,1.5) -- (4.7,1.5);
  \node[anchor=east, omDef!80!black, align=right] at (0.9,1.9) {बढ़ता हुआ\\ प्रोटीन};
  % incoming tRNA
  \draw[very thick, omMeth!80!black] (7.0,1.3) -- (7.0,2.1);
  \fill[omDef!70] (7.0,2.25) circle (4pt);
  \node[font=\scriptsize, omMeth!80!black, anchor=west] at (7.15,1.35) {CCA};
  \node[anchor=west, align=left, font=\scriptsize] at (7.3,1.7) {अगला tRNA, अपने\\ ऐमीनो अम्ल के साथ, आता हुआ};
  \draw[->, thick] (7.0,1.1) -- (6.1,0.75);
  \draw[->, thick] (5.7,-0.9) -- (7.2,-0.9) node[right] {राइबोसोम चलता है};
  \node[anchor=north, font=\scriptsize] at (4.05,-0.45) {कोडॉन};
  \draw (2.85,-0.35) -- (2.85,-0.45) -- (3.95,-0.45) -- (3.95,-0.35);
\end{tikzpicture}

{\small \omterm{prop:g11:gene-expression:translation}{अनुवादन}। \omterm{def:g10:cells-common-unit:organelle}{राइबोसोम} दो \omterm{def:g11:gene-expression:code}{कोडॉन} थामे रहता है; और जिस tRNA का
प्रतिकोडॉन किसी \omterm{def:g11:gene-expression:code}{कोडॉन} से मेल खाता है वह अपना \omterm{def:g10:chemistry-of-life:families}{ऐमीनो अम्ल} लाता है, शृंखला
नए आए \omterm{def:g10:chemistry-of-life:families}{ऐमीनो अम्ल} पर चढ़ा दी जाती है, और \omterm{def:g10:cells-common-unit:organelle}{राइबोसोम} एक \omterm{def:g11:gene-expression:code}{कोडॉन} आगे बढ़ जाता है।
विराम \omterm{def:g11:gene-expression:code}{कोडॉन} पर वह शृंखला छोड़ देता है।}
\end{omfigure}

\begin{example}[रफ़्तार और उपज]\label{ex:g11:gene-expression:speed}
कोई \omterm{def:g10:cells-common-unit:organelle}{राइबोसोम} प्रति सेकंड 5 से 20 \omterm{def:g10:chemistry-of-life:families}{ऐमीनो अम्ल} जोड़ता है: यानी 300 \omterm{def:g10:chemistry-of-life:families}{ऐमीनो
अम्लों} का \omterm{def:g11:gene-expression:protein}{प्रोटीन} मिनट भर से भी कम लेता है। \omterm{def:g10:cells-common-unit:organelle}{राइबोसोम} किसी mRNA पर लगभग 80
\omterm{def:g10:universal-dna:nucleotide}{न्यूक्लियोटाइड} की दूरी पर एक-दूसरे के पीछे चलते हैं, इसलिए 900
\omterm{def:g10:universal-dna:nucleotide}{न्यूक्लियोटाइडों} का संदेश एक साथ दर्जन भर ढोता है; और कुछ घंटों के अपने
जीवन में अकेला एक mRNA \omterm{def:g11:gene-expression:protein}{प्रोटीन} के हज़ार अणु बना देता है।
\emph{E.~coli} की एक \omterm{def:g10:cells-common-unit:cell}{कोशिका} प्रति सेकंड क़रीब $\num{20000}$ \omterm{def:g11:gene-expression:protein}{प्रोटीन} अणु
बनाती है; और लाल रक्त \omterm{def:g10:cells-common-unit:cell}{कोशिका} की पूर्ववर्ती अपने जीवन भर में क़रीब
$\num{5e8}$ हीमोग्लोबिन अणु।
\end{example}

\section{एक जीन, कई संदेश}

\begin{proposition}[सुकेंद्रकियों में संदेश का परिपक्व होना]\label{prop:g11:gene-expression:splicing}
सुकेंद्रकी \omterm{def:g10:cells-common-unit:cell}{कोशिकाओं} में किसी \omterm{def:g10:universal-dna:gene}{जीन} का अनुक्रम \emph{इंट्रॉनों}\index{इंट्रॉन}
से टूटा रहता है, यानी ऐसे खंडों से जिनका \omterm{def:g11:gene-expression:transcription}{अनुलेखन} तो होता है पर जो \omterm{def:g11:gene-expression:transcription}{RNA} के
\omterm{def:g10:cells-common-unit:organelle}{केंद्रक} छोड़ने से पहले उसमें से काट दिए जाते हैं; और जो खंड रखे जाकर सिरे
से सिरा जोड़ दिए जाते हैं, यानी \emph{एक्सॉन}, वही परिपक्व mRNA बनाते हैं।
बहुत-से \omterm{def:g10:universal-dna:gene}{जीन} एक से अधिक तरीक़ों से काटे जा सकते हैं, जिनमें एक्सॉनों के
अलग-अलग समुच्चय रखे जाते हैं (\emph{वैकल्पिक स्प्लाइसिंग}), जिससे एक ही
\omterm{def:g10:universal-dna:gene}{जीन} कई अलग mRNA और कई मिलते-जुलते \omterm{def:g11:gene-expression:protein}{प्रोटीन} देता है: मनुष्य के $\num{20000}$
\omterm{def:g10:universal-dna:gene}{जीन} $\num{100000}$ से कहीं अधिक प्रोटीनों के कूट हैं।
\end{proposition}
\admitted

\begin{omfigure}
\begin{tikzpicture}[font=\small, >=Stealth]
  \tikzset{ex/.style={draw, fill=omDef!40, minimum height=0.45cm, minimum width=1.0cm, inner sep=1pt},
           intr/.style={draw, fill=gray!20, minimum height=0.45cm, minimum width=0.7cm, inner sep=1pt}}
  \node[anchor=east] at (-0.2,0) {जीन (DNA)};
  \node[ex] at (0.5,0) {1}; \node[intr] at (1.35,0) {}; \node[ex] at (2.2,0) {2}; \node[intr] at (3.05,0) {}; \node[ex] at (3.9,0) {3}; \node[intr] at (4.75,0) {}; \node[ex] at (5.6,0) {4};
  \node[anchor=west, font=\scriptsize] at (6.2,0) {एक्सॉन (नीले), इंट्रॉन (धूसर)};
  \draw[->, thick] (3.0,-0.35) -- (3.0,-0.95) node[midway, right] {अनुलेखन};
  \node[anchor=east] at (-0.2,-1.3) {पूर्व-mRNA};
  \node[ex, fill=omThm!30] at (0.5,-1.3) {1}; \node[intr] at (1.35,-1.3) {}; \node[ex, fill=omThm!30] at (2.2,-1.3) {2}; \node[intr] at (3.05,-1.3) {}; \node[ex, fill=omThm!30] at (3.9,-1.3) {3}; \node[intr] at (4.75,-1.3) {}; \node[ex, fill=omThm!30] at (5.6,-1.3) {4};
  \draw[->, thick] (2.0,-1.65) -- (1.0,-2.35) node[midway, left] {स्प्लाइसिंग};
  \draw[->, thick] (4.0,-1.65) -- (6.0,-2.35) node[midway, right] {वैकल्पिक स्प्लाइसिंग};
  \node[anchor=east] at (-0.6,-2.7) {mRNA A};
  \node[ex, fill=omThm!30] at (0.0,-2.7) {1}; \node[ex, fill=omThm!30] at (1.0,-2.7) {2}; \node[ex, fill=omThm!30] at (2.0,-2.7) {3}; \node[ex, fill=omThm!30] at (3.0,-2.7) {4};
  \node[anchor=west] at (3.6,-2.7) {प्रोटीन A};
  \node[anchor=east] at (5.4,-2.7) {};
  \node[ex, fill=omThm!30] at (6.0,-2.7) {1}; \node[ex, fill=omThm!30] at (7.0,-2.7) {2}; \node[ex, fill=omThm!30] at (8.0,-2.7) {4};
  \node[anchor=west] at (8.6,-2.7) {प्रोटीन B};
\end{tikzpicture}

{\small सुकेंद्रकी \omterm{def:g10:universal-dna:gene}{जीन} का पूरा \omterm{def:g11:gene-expression:transcription}{अनुलेखन} होता है, फिर \omterm{prop:g11:gene-expression:splicing}{इंट्रॉन} हटा दिए जाते
हैं। सारे एक्सॉन रखने से, या उनमें से एक छोड़ देने से, एक ही \omterm{def:g10:universal-dna:gene}{जीन} से दो mRNA
और दो \omterm{def:g11:gene-expression:protein}{प्रोटीन} बनते हैं।}
\end{omfigure}

\begin{remark}[जानकारी का बहाव]\label{rem:g11:gene-expression:flow}
\omterm{def:g10:universal-dna:information}{DNA} का \omterm{def:g11:gene-expression:transcription}{अनुलेखन} \omterm{def:g11:gene-expression:transcription}{RNA} में होता है, \omterm{def:g11:gene-expression:transcription}{RNA} का \omterm{prop:g11:gene-expression:translation}{अनुवादन} \omterm{def:g11:gene-expression:protein}{प्रोटीन} में, और \omterm{def:g11:gene-expression:protein}{प्रोटीन}
लक्षण बनाता है: यानी जानकारी एक ही दिशा में बहती है। कोई बदला हुआ \omterm{def:g11:gene-expression:protein}{प्रोटीन}
\omterm{def:g10:universal-dna:gene}{जीन} को दोबारा कभी नहीं लिखता, और इसीलिए जीवन भर का \omterm{def:g10:sport-and-health:training}{प्रशिक्षण} या धूप की जलन
विरासत में नहीं जाती, तथा इसीलिए अगली पीढ़ी को जो मिलता है उसे केवल \omterm{def:g10:universal-dna:information}{DNA} के
\omterm{def:g11:mutations:mutation}{उत्परिवर्तन} (\cref{ch:g11:mutations}) ही बदलते हैं। ये तीनों क़दम हर जीव
की हर \omterm{def:g10:cells-common-unit:cell}{कोशिका} में, उसी कूट में चलते हैं: यानी
\cref{ch:g10:universal-dna} वाली सार्वत्रिकता, अब आख़िरी शब्द तक।
\end{remark}

\section{अभ्यास}

\begin{exercise}[$\star$]\label{exo:g11:gene-expression:1}
\omterm{def:g10:universal-dna:information}{DNA} और \omterm{def:g11:gene-expression:transcription}{RNA} के बीच तीन फ़र्क़ बताइए।
\end{exercise}

\begin{exercise}[$\star$]\label{exo:g11:gene-expression:2}
साँचा लड़ी \texttt{3'-TAC CGA AAA ATT-5'} का mRNA में \omterm{def:g11:gene-expression:transcription}{अनुलेखन} कीजिए।
\end{exercise}

\begin{exercise}[$\star$]\label{exo:g11:gene-expression:3}
कूट तालिका से mRNA \texttt{AUG GGC AAA UGU UAA} का \omterm{prop:g11:gene-expression:translation}{अनुवादन} कीजिए।
\end{exercise}

\begin{exercise}[$\star$]\label{exo:g11:gene-expression:4}
\omterm{prop:g11:gene-expression:translation}{अनुवादन} में \omterm{def:g10:cells-common-unit:organelle}{राइबोसोम} की और अंतरण \omterm{def:g11:gene-expression:transcription}{RNA} की क्या भूमिकाएँ हैं?
\end{exercise}

\begin{exercise}[$\star$]\label{exo:g11:gene-expression:5}
कूट को प्रति \omterm{def:g10:chemistry-of-life:families}{ऐमीनो अम्ल} कम से कम तीन \omterm{def:g10:universal-dna:nucleotide}{न्यूक्लियोटाइड} क्यों चाहिए?
\end{exercise}

\begin{exercise}[$\star\star$]\label{exo:g11:gene-expression:6}
किसी \omterm{def:g11:gene-expression:protein}{प्रोटीन} में 412 \omterm{def:g10:chemistry-of-life:families}{ऐमीनो अम्ल} हैं। विराम \omterm{def:g11:gene-expression:code}{कोडॉन} समेत उसके mRNA के कूट
वाले हिस्से की और उससे मेल खाते \omterm{def:g10:universal-dna:gene}{जीन} की, क्षारक-जोड़ों में, न्यूनतम लंबाई
क्या है?
\end{exercise}

\begin{exercise}[$\star\star$]\label{exo:g11:gene-expression:7}
कूट लड़ी \texttt{ATG AAA GGC TGG TAA} उत्परिवर्तित होकर
\texttt{ATG AAA GGC TGA TAA} बन जाती है। दोनों \omterm{def:g11:gene-expression:protein}{प्रोटीन} बताइए और
\omterm{def:g11:mutations:mutation}{उत्परिवर्तन} का वर्गीकरण कीजिए।
\end{exercise}

\begin{exercise}[$\star\star$]\label{exo:g11:gene-expression:8}
वही शुरुआती अनुक्रम \texttt{ATG AAG GGC TGG TAA} और
\texttt{ATG AAA GGG CTG GTA A} में बदलता है। हर एक का वर्गीकरण कीजिए और
दोनों \omterm{def:g11:gene-expression:protein}{प्रोटीन} बताइए।
\end{exercise}

\begin{exercise}[$\star\star$]\label{exo:g11:gene-expression:9}
तालिका की मदद से समझाइए कि किसी \omterm{def:g11:gene-expression:code}{कोडॉन} के तीसरे स्थान पर प्रतिस्थापन अक्सर
मूक क्यों रहता है।
\end{exercise}

\begin{exercise}[$\star\star$]\label{exo:g11:gene-expression:10}
पॉली-UC वाला \omterm{def:g11:gene-expression:transcription}{RNA} (\texttt{UCUCUCUC}\dots) ऐसा \omterm{def:g11:gene-expression:protein}{प्रोटीन} देता है जिसमें
सेरीन और ल्यूसीन बारी-बारी से आते हैं। दिखाइए कि यह तीन अक्षरों वाले कूट
से संगत है और उससे दो \omterm{def:g11:gene-expression:code}{कोडॉन} तय कीजिए।
\end{exercise}

\begin{exercise}[$\star\star$]\label{exo:g11:gene-expression:11}
\num{30000} क्षारक-जोड़ों वाला कोई मानव \omterm{def:g10:universal-dna:gene}{जीन} 500 \omterm{def:g10:chemistry-of-life:families}{ऐमीनो अम्लों} का \omterm{def:g11:gene-expression:protein}{प्रोटीन}
बनाता है। इस बेमेल को समझाइए।
\end{exercise}

\begin{exercise}[$\star\star\star$]\label{exo:g11:gene-expression:12}
कोई दवा जीवाणुओं के \omterm{def:g10:cells-common-unit:organelle}{राइबोसोम} रोक देती है पर मनुष्य के नहीं। समझाइए कि वह
प्रतिजैविक क्यों हो सकती है, और उसका होना दोनों समूहों के \omterm{def:g10:cells-common-unit:organelle}{राइबोसोमों} के
बारे में क्या बताता है।
\end{exercise}

\begin{exercise}[$\star\star\star$]\label{exo:g11:gene-expression:13}
इस अध्याय की शब्दावली में समझाइए कि \cref{ch:g10:universal-dna} वाला
जेलीफ़िश का \omterm{def:g10:universal-dna:gene}{जीन} कोई चूहा कैसे पढ़ सका, और चूहा कोई और कूट इस्तेमाल करता तो
क्या होता।
\end{exercise}

\begin{exercise}[$\star\star\star$]\label{exo:g11:gene-expression:14}
कोई \omterm{def:g11:mutations:mutation}{उत्परिवर्तन} उस tRNA को बदल देता है जिसका प्रतिकोडॉन UAG पढ़ता है,
ताकि वह रोकने के बजाय कोई \omterm{def:g10:chemistry-of-life:families}{ऐमीनो अम्ल} ढो लाए। \omterm{def:g10:cells-common-unit:cell}{कोशिका} के प्रोटीनों पर आम तौर
पर, और समय से पहले UAG ढोते किसी उत्परिवर्ती \omterm{def:g10:universal-dna:gene}{जीन} पर, उसके असर का अनुमान
लगाइए।
\end{exercise}

\begin{exercise}[$\star\star\star$]\label{exo:g11:gene-expression:15}
चर्चा कीजिए कि किसी \omterm{def:g10:universal-dna:gene}{जीन} के आरंभ के पास का चौखट-खिसकाव उसके अंत के पास वाले
से लगभग हमेशा बुरा क्यों होता है, और कोई कुअर्थी \omterm{def:g11:mutations:mutation}{उत्परिवर्तन} हानिरहित से
लेकर जानलेवा तक कुछ भी क्यों हो सकता है।
\end{exercise}

\section{समस्या: एक जीन, चार तरह से पढ़ा गया}

\begin{problem}[{सप्ताहांत समस्या --- एक छोटा जीन अनुलिखित, अनुवादित, तीन
बार उत्परिवर्तित और समयबद्ध किया गया: DNA से ऐमीनो अम्लों तक, और वे साठ
सेकंड जो किसी प्रोटीन में लगते हैं}]
\label{pb:g11:gene-expression:1}
किसी छोटे \omterm{def:g10:universal-dna:gene}{जीन} की कूट लड़ी यह पढ़ी जाती है:
\begin{center}
\texttt{ATG GCT TGG AAA CCC GAA TAC GGT CAT TTC TGA}
\end{center}

\medskip
\noindent\textbf{भाग I --- \omterm{def:g10:universal-dna:gene}{जीन} पढ़ना।}
\begin{enumerate}
  \item साँचा लड़ी लिखिए।
  \item \omterm{def:g10:universal-dna:gene}{जीन} से अनुलिखित mRNA लिखिए।
  \item उसका कोडॉन-दर-कोडॉन \omterm{prop:g11:gene-expression:translation}{अनुवादन} कीजिए और \omterm{def:g11:gene-expression:protein}{प्रोटीन} के \omterm{def:g10:chemistry-of-life:families}{ऐमीनो अम्लों} का
        अनुक्रम बताइए।
  \item उस \omterm{def:g11:gene-expression:protein}{प्रोटीन} में कितने \omterm{def:g10:chemistry-of-life:families}{ऐमीनो अम्ल} हैं, और उन्हें तय करने में, विराम
        समेत, mRNA के कितने \omterm{def:g10:universal-dna:nucleotide}{न्यूक्लियोटाइड} लगते हैं?
  \item पढ़ना किस \omterm{def:g11:gene-expression:code}{कोडॉन} से शुरू हुआ, और बाक़ी सबके लिए वाचन चौखट किससे तय
        होती है?
\end{enumerate}

\medskip
\noindent\textbf{भाग II --- तीन उत्परिवर्ती।}
\begin{enumerate}[resume]
  \item उत्परिवर्ती 1: \texttt{ATG GCT TGG AAA CCG GAA TAC GGT CAT TTC TGA}।
        \omterm{def:g11:gene-expression:protein}{प्रोटीन} बताइए और \omterm{def:g11:mutations:mutation}{उत्परिवर्तन} का वर्गीकरण कीजिए।
  \item उत्परिवर्ती 2: \texttt{ATG GCT TGG AAA CCC GAA TAC GGT CAT TTC TGA}
        जिसमें नौवाँ \omterm{def:g11:gene-expression:code}{कोडॉन} \texttt{CAT} बदलकर \texttt{CGT} हो गया है।
        \omterm{def:g11:gene-expression:protein}{प्रोटीन} बताइए और वर्गीकरण कीजिए।
  \item उत्परिवर्ती 3: \texttt{ATG GCT TGA AAA CCC GAA TAC GGT CAT TTC TGA}।
        \omterm{def:g11:gene-expression:protein}{प्रोटीन} बताइए और वर्गीकरण कीजिए। कौन-से अकेले प्रतिस्थापन ने उसे
        बनाया?
  \item उत्परिवर्ती 4: छठे \omterm{def:g10:universal-dna:nucleotide}{न्यूक्लियोटाइड} के बाद एक T डाला गया है:
        \texttt{ATG GCT TTG GAA ACC CGA ATA CGG TCA TTT CTG A}। \omterm{def:g11:gene-expression:protein}{प्रोटीन}
        बताइए और वर्गीकरण कीजिए।
  \item \omterm{def:g11:gene-expression:protein}{प्रोटीन} के लिए चारों उत्परिवर्तियों को सबसे कम से सबसे अधिक
        विघटनकारी के क्रम में रखिए और कारण बताइए।
\end{enumerate}

\medskip
\noindent\textbf{भाग III --- कूट ऐसा क्यों है।}
\begin{enumerate}[resume]
  \item ल्यूसीन के कितने \omterm{def:g11:gene-expression:code}{कोडॉन} हैं? सेरीन के? ट्रिप्टोफ़ैन के? किसी
        यादृच्छिक प्रतिस्थापन से अपने \omterm{def:g11:gene-expression:code}{कोडॉन} में बदल जाने की सबसे अधिक
        संभावना किस \omterm{def:g10:chemistry-of-life:families}{ऐमीनो अम्ल} की है, और सबसे कम किसकी?
  \item किसी ल्यूसीन \omterm{def:g11:gene-expression:code}{कोडॉन} \texttt{CUx} के तीसरे स्थान पर होने वाले सारे
        प्रतिस्थापनों में से मूक रहने वालों का अंश निकालिए।
  \item समझाइए कि दो अक्षरों वाला कूट क्यों नहीं चल सकता था, और चार
        अक्षरों वाला क्यों नहीं चाहिए।
  \item पॉली-U पॉली-फेनिलऐलानिन देता है और पॉली-A पॉली-लाइसीन। ये दोनों
        प्रयोग तालिका की कौन-सी प्रविष्टियाँ स्थापित करते हैं?
  \item कूट जीवाणु में और मनुष्य में एक ही है। \omterm{prop:g10:universal-dna:universal}{पारजीनन} के लिए इसका नतीजा
        बताइए, और नातेदारी के लिए भी।
\end{enumerate}

\medskip
\noindent\textbf{भाग IV --- कारख़ाने की समय-सारणी।}
कोई \omterm{def:g10:cells-common-unit:organelle}{राइबोसोम} प्रति सेकंड 10 \omterm{def:g10:chemistry-of-life:families}{ऐमीनो अम्ल} जोड़ता है और \omterm{def:g10:cells-common-unit:organelle}{राइबोसोम} किसी mRNA के
साथ-साथ 80 \omterm{def:g10:universal-dna:nucleotide}{न्यूक्लियोटाइड} की दूरी पर एक-दूसरे के पीछे चलते हैं; और हर mRNA
2 घंटे जीता है।
\begin{enumerate}[resume]
  \item इस \omterm{def:g10:universal-dna:gene}{जीन} के \omterm{def:g11:gene-expression:protein}{प्रोटीन} का \omterm{prop:g11:gene-expression:translation}{अनुवादन} करने में एक \omterm{def:g10:cells-common-unit:organelle}{राइबोसोम} को कितना समय
        लगता है? और 600 \omterm{def:g10:chemistry-of-life:families}{ऐमीनो अम्लों} के \omterm{def:g11:gene-expression:protein}{प्रोटीन} में?
  \item 1800 \omterm{def:g10:universal-dna:nucleotide}{न्यूक्लियोटाइडों} वाले किसी mRNA को एक साथ कितने \omterm{def:g10:cells-common-unit:organelle}{राइबोसोम} पढ़
        सकते हैं?
  \item यदि हर 8 सेकंड में एक नया \omterm{def:g10:cells-common-unit:organelle}{राइबोसोम} शुरू हो, तो एक mRNA अपने जीवन
        भर में उस \omterm{def:g11:gene-expression:protein}{प्रोटीन} की कितनी प्रतियाँ देता है?
  \item किसी जीवाणु को 10 मिनट में किसी \omterm{def:g10:cell-metabolism:metabolism}{एंजाइम} की \num{2000} प्रतियाँ
        चाहिए। एक साथ अनुलिखित उस \omterm{def:g10:universal-dna:gene}{जीन} के कितने mRNA चाहिए?
  \item परिणाम लिखिए: उस \omterm{def:g10:universal-dna:gene}{जीन} के कूट वाला \omterm{def:g11:gene-expression:protein}{प्रोटीन}, चारों में से वह
        \omterm{def:g11:mutations:mutation}{उत्परिवर्तन} जिसने उसे मिटा दिया, और एक \omterm{def:g10:cells-common-unit:organelle}{राइबोसोम} को उसे बनाने में
        लगने वाला समय।
\end{enumerate}
\end{problem}
